\chapter{Algoritmicky těžké problémy}\label{chap:tezke-problemy}

Mnoho problémů lze v~informatice řešit pomocí polynomiálních algoritmů,~tzn. algoritmů běžících v~čase $\bigO{n^k}$ pro nějaké pevné $k$. Například prvky v~poli můžeme řadit různými algoritmy; jednoduchý \emph{BubbleSort} má v~nejhorším případě časovou složitost $\bigO{n^2}$,~kde $n$ je počet prvků pole. Také nejkratší cesty v~grafu dokážeme hledat polynomiálně,~třeba Dijkstrovým algoritmem v~čase $\bigO{n^2+m}$ při použití pole,~kde $n$ je počet vrcholů a~$m$ počet hran.\footnote{U~jednoduchého neorientovaného grafu je $m\leqslant n(n-1)/2$,~u~jednoduchého orientovaného grafu $m\leqslant n(n-1)$. V~obou případech je tedy $m$ polynomiálně omezeno pomocí $n$.} S~jistou rezervou lze říci,~že polynomiální algoritmy bývají prakticky použitelné.

Bohužel ne vždy je situace tak příznivá. Lze narazit na problémy,~pro které žádný polynomiální algoritmus neznáme,~a~u~některých dokonce umíme dokázat,~že je v~polynomiálním čase řešit nelze. Dobrým příkladem úlohy s~prokazatelně exponenciálně dlouhým řešením jsou \emph{Hanojské věže}: přesunutí $n$ disků vyžaduje nejméně $2^n-1$ tahů.

Nás budou zajímat ty nejdůležitější problémy z~této oblasti,~protože mezi nimi lze nalézt zajímavé vztahy,~a~posléze si uděláme menší ochutnávku z~problematiky \P{} a~\NP.

\section{Problém splnitelnosti}\label{sec:general_sat}

Začneme zdánlivě možná jednoduchým problémem,~který však na poli informatiky hraje dosti důležitou roli. Předpokládejme,~že máme zadanou nějakou logickou formuli $\varphi$ o~libovolném počtu logických proměnných,~označme např. $x_1,x_2,\dots,x_n$. Naší otázkou je,~jestli je možné hodnoty proměnných $x_1,x_2,\dots,x_n$ nastavit tak (tj. dosadit hodnoty 0~a~1),~aby byla formule splněna\footnote{Může se na první pohled zdát,~že se jedná o~dosti teoretickou záležitost,~nicméně později uvidíme,~že mnoho jiných problémů má se \ensuremath{\SAT} silnou spojitost.},~tj. $\varphi(\dots)=1$. Takovému ohodnocení budeme říkat \emph{splňující}.

\problem{Splnitelnost obecné logické formule}{Logická formule $\varphi$.}{1,~pokud je $\varphi$ splnitelná,~jinak 0.}

Nejdříve si zopakujeme logické operace a~spojky. Mezi ně řadíme $\neg,~\land,~\lor,~\implies,~\iff$.
\begin{itemize}
    \item \textbf{Negace $\neg x$.} \emph{Převrací logickou hodnotu $x$.}
    \item \textbf{Konjunkce $a\land b$.} \emph{Pravdivá,~jsou-li obě hodnoty operandů $a,b$ pravdivé.}
    \item \textbf{Disjunkce $a\lor b$.} \emph{Pravdivá,~je-li alespoň jedna z~hodnot operandů $a,b$ pravdivá.}
    \item \textbf{Implikace $a\implies b$.} \emph{Je-li předpoklad $a$ splněn,~je splněn i~závěr $b$. Tj. implikace je nepravdivá,~pokud platí $a$ a~neplatí $b$.}
    \item \textbf{Ekvivalence $a\iff b$.} \emph{$a$ platí právě tehdy,~když platí $b$. Tj. ekvivalence je pravdivá,~jsou-li hodnoty operandů $a,b$ současně 0~nebo 1.}
\end{itemize}
\begin{table}[ht]\label{table:logicke_operace}
    \centering
    \begin{tabular}{|c|c|c|c|c|c|c|}
    \hline
    $a$ & $b$ & $\neg a$ & $a\land b$ & $a\lor b$ & $a\implies b$ & $a\iff b$ \\ \hline
    0   & 0   & 1        & 0          & 0         & 1             & 1         \\ \hline
    0   & 1   & 1        & 0          & 1         & 1             & 0         \\ \hline
    1   & 0   & 0        & 0          & 1         & 0             & 0         \\ \hline
    1   & 1   & 0        & 1          & 1         & 1             & 1         \\ \hline
    \end{tabular}
\end{table}
Zde konstatujme,~že ve všech příkladech bude mít negace $\neg$ vždy \emph{nejvyšší prioritu} oproti ostatním operacím. Prioritu ostatních operací budeme stanovovat pomocí závorek. 
\begin{example}\label{ex:sat_formule}
    \begin{itemize}
        \item $\varphi(x)=x\land\neg x$ \dots \textbf{ Není splnitelná},~pro $x=0$ i~$x=1$ je $\varphi=0$.
        \item $\psi(x)=x\lor\neg x$ \dots \textbf{ Je splnitelná},~a~to jak pro $x=0$,~tak $x=1$.
        \item $\chi(a,b,c)=\big((a\land b)\lor \neg c\big)\implies (\neg a\land\neg c)$ \dots \textbf{Je splnitelná},~např. pro $a=0$,~$b=0$ a~$c=0$.
    \end{itemize}
\end{example}
V~příkladu \ref{ex:sat_formule} výše se jedná o~poměrně jednoduché instance,~kdy jsme schopni vidět hned,~zda je formule splnitelná. Pokud bychom však uvážili nějakou složitější formuli,~problém se již značně zkomplikuje.
\importantbox{Naivní postup pro formuli o~$n$ proměnných prozkoumá až $2^n$ možných ohodnocení. Algoritmus zkoušející všechny možnosti je tak \emph{exponenciální}.}
Pro \ensuremath{\SAT} není dodnes známý žádný algoritmus,~který by jej uměl řešit pro libovolnou logickou formuli v~polynomiálním čase.

\subsection{Konjunktivní normální forma}\label{subsec:cnf}

Zkusíme se podívat na formule ve speciálním tvaru,~a~to tzv. \emph{konjunktivní normální formě},~neboli zkráceně \emph{CNF}.\footnote{Z~angl. \emph{conjunctive normal form}.}

Formule v~CNF se skládá z~\textbf{klauzulí},~které obsahují \textbf{literály}. Literálem je buď proměnná,~nebo její negace.
\begin{itemize}
    \item Klauzule jsou ve formuli odděleny \emph{logickou spojkou $\land$}
    \item a~v~každé klauzuli jsou literály odděleny \emph{logickou spojkou $\lor$}.
\end{itemize}
\[\varphi(\dots)=\underbrace{(\cdots\lor\cdots\lor\cdots)}_{\text{klauzule}}\land(\cdots\lor \overbrace{x}^{\text{literál}}\lor\cdots)\land(\cdots\lor\overbrace{\neg x}^{\text{literál}}\lor\cdots)\land(\cdots\lor \underbrace{y}_{\text{literál}}\lor\cdots)\land\cdots\]
\begin{example}\label{ex:cnf_formule}
    \begin{itemize}
        \item \(\varphi(x,y,z)=\neg x\land (y\lor z)\) \dots \textbf{Je v~CNF.} Formule $\varphi$ obsahuje klauzule $\neg x$ (klauzule o~jednom literálu) a~$y\lor z$.
        \item \(\psi(a,b)=a\land b\) \dots \textbf{Je v~CNF.}
        \item \(\chi(a,b,c,d)=a\land \big(b\lor (c\land d)\big)\) \dots \textbf{Není v~CNF},~protože výraz $c\land d$ není literál.
    \end{itemize}
    Poslední formuli lze však převést do CNF: $a\land (b\lor c)\land (b\lor d)$.
\end{example}
Z~příkladu \ref{ex:cnf_formule} se nabízí otázka,~zda je možné \emph{libovolnou formuli} převést do CNF. Existuje pro každou formuli ekvivalentní\footnote{Formule $\varphi$ a~$\psi$ jsou ekvivalentní,~pokud mají při každém ohodnocení společných proměnných stejnou pravdivostní hodnotu.} formule v~CNF? Ano.
\begin{theorem}
    Pro každou logickou formuli $\psi$ existuje ekvivalentní formule $\psi^\prime$ v~CNF.
\end{theorem}
Toto tvrzení lze,~pochopitelně,~dokázat,~ale zde se tím zabývat nebudeme a~zkrátka jej přijmeme jako fakt. Podstatná věc,~která z~toho plyne,~je,~že pokud je nějaká formule $\psi$ splnitelná,~pak je splnitelná i~jí ekvivalentní formule $\psi^\prime$ v~CNF a~naopak pokud $\psi$ není splnitelná,~není splnitelná ani $\psi^\prime$. Symbolicky
\[\psi\,\text{je splnitelná}\iff\psi^\prime\,\text{je splnitelná}.\]

\subsection{Převod do CNF pomocí tabulky}\label{subsec:cnf_prevod_tabulka}

Podívejme se na způsob,~jak lze převést libovolnou logickou formuli do CNF.

Máme obecnou formuli $\varphi$ s~proměnnými $x_1,x_2,\dots,x_n$,~pro kterou budeme konstruovat ekvivalentní formuli $\varphi^\prime$ v~CNF. Vycházíme z~tabulky pravdivostních hodnot,~přičemž pro každé \textbf{nepravdivé} ohodnocení vytvoříme \emph{klauzuli} s~$n$ literály,~kde obecně $i$-tý literál je
\begin{itemize}
    \item $x_i$,~pokud v~daném ohodnocení je $x_i=0$,
    \item jinak $\neg x_i$ (tj. když $x_i=1$).
\end{itemize}
Podívejme se na úvod na příklad \ref{ex:prevod_cnf_1} níže.
\begin{example}\label{ex:prevod_cnf_1}
    Máme formuli $\psi(x,y,z)=(x\implies y)\implies z$,~pro kterou chceme najít ekvivalentní formuli $\psi^\prime$ v~CNF. Nejprve si tedy sestavíme tabulku pravdivostních hodnot.
    \begin{table}[ht]
        \centering
        \begin{tabular}{|c|c|c|c|c|}
        \hline
        $x$ & $y$ & $z$ & $x\implies y$ & $(x\implies y)\implies z$ \\ \hline
        0   & 0   & 0   & 1             & 0                         \\ \hline
        0   & 0   & 1   & 1             & 1                         \\ \hline
        0   & 1   & 0   & 1             & 0                         \\ \hline
        0   & 1   & 1   & 1             & 1                         \\ \hline
        1   & 0   & 0   & 0             & 1                         \\ \hline
        1   & 0   & 1   & 0             & 1                         \\ \hline
        1   & 1   & 0   & 1             & 0                         \\ \hline
        1   & 1   & 1   & 1             & 1                         \\ \hline
        \end{tabular}
    \end{table}
    Z~tabulky můžeme vidět,~že formule $\psi$ je nepravdivá pro
    \begin{itemize}
        \item $x=0,y=0,z=0$,
        \item $x=0,y=1,z=0$
        \item a~$x=1,y=1,z=0$.
    \end{itemize}
    Tedy celkově sestavíme 3~klauzule:
    \begin{itemize}
        \item \((\overbrace{x}^{x=0}\lor \underbrace{y}_{y=0}\lor \overbrace{z}^{z=0})\),
        \item \((x\lor\underbrace{\neg y}_{y=1}\lor z)\),
        \item \((\neg x\lor\neg y\lor z)\).
    \end{itemize}
    Výsledná formule v~CNF bude mít tvar:
    \[\psi^\prime(x,y,z)=(x\lor y\lor z)\land (x\lor \neg y\lor z)\land (\neg x\lor \neg y\lor z).\]
    Čtenář se sám může přesvědčit,~že formule $\psi$ a~$\psi^\prime$ pro libovolné ohodnocení mají stejnou pravdivostní hodnotu.
\end{example}
Nabízí se otázka: \emph{\uv{Proč tento postup funguje?}} Pokud formule v~CNF není při určitém ohodnocení pravdivá,~alespoň jedna její klauzule není pravdivá. Klauzule obsahující všechny proměnné je nepravdivá právě při jediném ohodnocení: všechny její literály musí být nepravdivé. Pro každé ohodnocení,~při němž je původní formule nepravdivá,~tedy sestrojíme klauzuli,~která zakáže právě toto ohodnocení.
\begin{example}
    \(\chi(x_1,x_2,x_3,x_4)=(x_1\iff x_2)\lor \neg(x_3\iff x_4)\).
    \begin{table}[ht]
        \centering
        \begin{tabular}{|c|c|c|c|c|c|c|}
        \hline
        $x_1$ & $x_2$ & $x_3$ & $x_4$ & $x_1\iff x_2$ & $\neg(x_3\iff x_4)$ & $(x_1\iff x_2)\lor\neg(x_3\iff x_4)$ \\ \hline
        0     & 0     & 0     & 0     & 1             & 0                                                 & 1                                    \\ \hline
        0     & 0     & 0     & 1     & 1             & 1                                                 & 1                                    \\ \hline
        0     & 0     & 1     & 0     & 1             & 1                                                 & 1                                    \\ \hline
        0     & 0     & 1     & 1     & 1             & 0                                                 & 1                                    \\ \hline
        0     & 1     & 0     & 0     & 0             & 0                                                 & 0                                    \\ \hline
        0     & 1     & 0     & 1     & 0             & 1                                                 & 1                                    \\ \hline
        0     & 1     & 1     & 0     & 0             & 1                                                 & 1                                    \\ \hline
        0     & 1     & 1     & 1     & 0             & 0                                                 & 0                                    \\ \hline
        1     & 0     & 0     & 0     & 0             & 0                                                 & 0                                    \\ \hline
        1     & 0     & 0     & 1     & 0             & 1                                                 & 1                                    \\ \hline
        1     & 0     & 1     & 0     & 0             & 1                                                 & 1                                    \\ \hline
        1     & 0     & 1     & 1     & 0             & 0                                                 & 0                                    \\ \hline
        1     & 1     & 0     & 0     & 1             & 0                                                 & 1                                    \\ \hline
        1     & 1     & 0     & 1     & 1             & 1                                                 & 1                                    \\ \hline
        1     & 1     & 1     & 0     & 1             & 1                                                 & 1                                    \\ \hline
        1     & 1     & 1     & 1     & 1             & 0                                                 & 1                                    \\ \hline
        \end{tabular}
    \end{table}
    \[\chi^\prime(x_1,x_2,x_3,x_4)=(x_1\lor\neg x_2\lor x_3\lor x_4)\land(x_1\lor\neg x_2\lor\neg x_3\lor\neg x_4)\land(\neg x_1\lor x_2\lor x_3\lor x_4)\land(\neg x_1\lor x_2\lor\neg x_3\lor\neg x_4).\]
    \importantbox{\textbf{Problém:} Generování tabulky (resp. procházení všech možných ohodnocení) nás přivádí zpět k~původnímu problému,~a~to sice \emph{exponenciální časové složitosti}.}
    Tento postup nelze obecně zachránit pouhou chytřejší implementací: pokud trváme na ekvivalentní CNF bez nových proměnných,~může být výsledná formule exponenciálně delší. Například při roznásobení formule
    \[(x_1\land y_1)\lor(x_2\land y_2)\lor\dots\lor(x_n\land y_n)\]
    vzniká pro každou volbu jednoho literálu z~každé dvojice samostatná klauzule,~tedy celkem $2^n$ klauzulí.
\end{example}

\subsection{Tseitinova transformace}\label{subsec:cnf_tseitin}

Značně rychlejší převod do CNF navrhl sovětský matematik \emph{Grigorij S.~Cejtin}\footnote{Jeho příjmení se v~anglických textech obvykle přepisuje jako \emph{Tseitin}; odtud pochází ustálený název transformace.}. Předchozí postup trval na tom,~že ve výsledku ponecháme pouze původní proměnné. Tseitinova transformace zavede pomocné proměnné a~získá formuli,~která sice obecně není s~původní formulí ekvivalentní nad rozšířenou množinou proměnných,~ale je s~ní \emph{stejně splnitelná}.

Postup lze rozdělit do čtyř kroků.
\begin{enumerate}
    \item Formuli rozložíme na podformule odpovídající uzlům jejího syntaktického stromu.
    \item Pro každou složenou podformuli zavedeme novou proměnnou.
    \item Pro každou novou proměnnou přidáme klauzule,~které vynutí,~aby měla stejnou hodnotu jako příslušná podformule.
    \item Nakonec přidáme jednotkovou klauzuli tvořenou proměnnou zastupující celou původní formuli; tím požadujeme,~aby její hodnota byla $1$.
\end{enumerate}

Potřebné klauzule jsou krátké a~lze je zapsat přímo. Pro literály nebo pomocné proměnné $a,b$ platí
\begin{align}
y\iff\neg a
    &\quad\leadsto\quad(\neg y\lor\neg a)\land(y\lor a),\label{eq:tseitin-neg}\\
y\iff(a\land b)
    &\quad\leadsto\quad(\neg y\lor a)\land(\neg y\lor b)\land(y\lor\neg a\lor\neg b),\label{eq:tseitin-and}\\
y\iff(a\lor b)
    &\quad\leadsto\quad(y\lor\neg a)\land(y\lor\neg b)\land(\neg y\lor a\lor b).\label{eq:tseitin-or}
\end{align}
Implikaci nejprve přepíšeme pomocí $a\implies b\equiv\neg a\lor b$. Každý výskyt spojky tak přidá nejvýše tři klauzule konstantní délky.

Pojďme se podívat na konkrétní příklad.
\begin{example}
    Máme formuli $\alpha(p,q,r)=((p\lor q)\land r)\implies \neg p$. Chceme pro ni sestrojit Tseitinovou transformací novou formuli $\beta$.
    
    Formule $\alpha$ obsahuje tyto dílčí podformule:
    \begin{itemize}
        \item $\neg p$,
        \item $p\lor q$,
        \item $(p\lor q)\land r$,
        \item a~$((p\lor q)\land r)\implies \neg p$ (původní formule $\alpha$).
    \end{itemize}
    Pro ně zavedeme proměnné $y_1,y_2,y_3$ a~$y_4$. Implikaci nejprve přepíšeme jako disjunkci:
    \begin{itemize}
        \item $y_4\iff\neg p$,
        \item $y_3\iff(p\lor q)$,
        \item $y_2\iff(y_3\land r)$
        \item a~$y_1\iff(\neg y_2\lor y_4)$.
    \end{itemize}
    Pomocí vztahů \eqref{eq:tseitin-neg}--\eqref{eq:tseitin-or} dostaneme formuli v~CNF:
    \begin{align*}
        \beta={}&y_1\\
        &\land(y_1\lor y_2)\land(y_1\lor\neg y_4)
            \land(\neg y_1\lor\neg y_2\lor y_4)\\
        &\land(\neg y_2\lor y_3)\land(\neg y_2\lor r)
            \land(y_2\lor\neg y_3\lor\neg r)\\
        &\land(y_3\lor\neg p)\land(y_3\lor\neg q)
            \land(\neg y_3\lor p\lor q)\\
        &\land(\neg y_4\lor\neg p)\land(y_4\lor p).
    \end{align*}

    Zkusme ohodnocení $p=1,q=0$ a~$r=0$. Původní formule je pravdivá:
    \[\alpha(1,0,0)=((1\lor 0)\land 0)\implies\neg 1=0\implies 0=1.\]
    Pomocné proměnné musí nabývat hodnot
    \begin{itemize}
        \item $y_4=0$,~protože $\neg p=0$,
        \item $y_3=1$,~protože $p\lor q=1$,
        \item $y_2=0$,~protože $y_3\land r=0$,
        \item a~$y_1=1$,~protože $\neg y_2\lor y_4=1$.
    \end{itemize}
    Po dosazení jsou všechny klauzule formule $\beta$ splněny.
\end{example}

Je-li původní formule splnitelná,~nastavíme každou pomocnou proměnnou podle hodnoty její podformule a~splníme všechny nové klauzule. Je-li naopak výsledná formule splnitelná,~klauzule vynucují správné hodnoty všech pomocných proměnných a~jednotková klauzule pro kořen říká,~že původní formule musí být pravdivá. Obě formule jsou tedy stejně splnitelné.
\tipbox{Každá spojka původní formule přidá jednu pomocnou proměnnou a~nejvýše tři klauzule. Velikost výsledné formule i~doba konstrukce jsou proto lineární ve velikosti původní formule.}

\section{Problém \ensuremath{\SAT}}\label{sec:sat}

V~podsekci \ref{subsec:cnf} jsme viděli,~že převod obecné formule do CNF \uv{hrubou silou} je problematický. Nová formule se totiž může až exponenciálně prodloužit. V~podsekci \ref{subsec:cnf_tseitin} jsme si proto ukázali Tseitinovu transformaci,~která pomocí nových proměnných sestrojí stejně splnitelnou CNF efektivně. Zkusme si tedy původní problém zjednodušit. Předpokládejme nyní,~že formule na vstupu jsou již v~CNF. Tento problém je v~informatice velmi význačný a~označuje se zkratkou \ensuremath{\SAT} (z~angl. \emph{satisfiability},~neboli splnitelnost).

\problem{Splnitelnost logické formule v~CNF (\ensuremath{\SAT})}{Logická formule $\varphi$ v~CNF.}{1,~pokud je $\varphi$ splnitelná,~jinak 0.}

Nyní se zaměříme na samotnou splnitelnost vstupní formule. V~praxi se používají tzv. \emph{\ensuremath{\SAT} solvery},~které rozhodují,~zda je formule v~CNF splnitelná. Ukážeme si klasický algoritmus \emph{DPLL},~jehož název připomíná příjmení Martina Davise,~Hilaryho Putnama,~George Logemanna a~Donalda Lovelanda. Jeho základní podoba pochází z~počátku 60.~let 20.~století.

\subsection{Algoritmus DPLL}\label{subsec:dpll}

Tento algoritmus je založený na rekurzivním postupu. Na vstupu algoritmu je libovolná formule $\varphi$ v~CNF,~kterou algoritmus postupně vyhodnocuje. Využívá dvojici procedur -- \textbf{Unit propagation} a~\textbf{Pure literal elimination}.

\subsubsection*{Unit propagation}

Tato procedura hledá v~dané formuli $\varphi$ tzv. \emph{jednotkovou klauzuli}. To je klauzule,~která obsahuje \textbf{pouze jeden literál}. Například formule $\omega$ níže obsahuje jednotkovou klauzuli.
\[\omega(x,y)=(x\lor\neg y)\land \underbrace{\neg y}_{\text{jednotková kl.}}\land (\neg x\lor y).\]
Pokud se taková klauzule nachází ve formuli,~lze toho využít,~neboť pro danou proměnnou existuje právě jedno ohodnocení,~které tuto klauzuli splní. Obecně pokud formule $\varphi$ obsahuje jednotkovou klauzuli s~proměnnou $x$,~pak mohou nastat dva případy. Označme $\varphi^\prime$ zbytek klauzulí ve formuli $\varphi$.
\begin{enumerate}
    \item Pokud má formule tvar $\varphi=x\land\varphi^\prime$,~pak nutně $x=1$.
    \item Jinak pokud má formule tvar $\varphi=\neg x\land\varphi^\prime$,~pak musí být $x=0$.
\end{enumerate}
Fakt,~že jsme schopni hned rozhodnout,~jaké ohodnocení musí mít daná proměnná,~je v~tomto případě hodně ku prospěchu,~neboť formuli můžeme zjednodušit. Protože ohodnocení dané jednotkové klauzule již známe,~můžeme ji z~formule $\varphi$ odstranit. Zároveň můžeme \textbf{odstranit všechny klauzule},~které se díky proměnné $x$ vyhodnotí na $1$,~neboť již nemohou nijak ovlivnit logickou hodnotu zbylých klauzulí.

Např. pokud formule obsahuje literál $x$ v~jednotkové klauzuli,~pak hodnota proměnné musí být $1$. Tzn. klauzule obsahující $x$ se vyhodnotí též na $1$. tj. 
\[(\dots\lor x\lor\dots)=1.\]
Naopak pokud obsahuje literál $\neg x$,~pak nutně $x=0$ a~tedy
\[(\dots\lor \neg x\lor\dots)=1.\]
Dále ještě odstraníme výskyty literálu $x$,~resp. $\neg x$ v~\emph{opačné polaritě},~tedy negaci původního literálu. Tzn. pokud literál v~jednotkové klauzuli byl $x$,~pak opačná polarita literálu je $\neg x$ a~naopak. Důvod,~proč jej můžeme odstranit,~je ten,~že v~rámci kterékoliv jiné klauzule se daný literál v~opačné polaritě vyhodnotí vždy na $0$. Protože je však spojen s~ostatními literály pomocí $\lor$,~nemůže nijak ovlivnit výslednou logickou hodnotu klauzule a~tudíž jej můžeme odstranit. Tedy např. pokud $x=1$,~pak můžeme klauzule tvaru
\[(\dots\lor y \lor \underbrace{\neg x}_{\text{odstraníme}} \lor z~\lor \dots)\]
zjednodušit na
\[(\dots\lor y \lor z~\lor \dots).\]
Podobně pro $x=0$:
\[(\dots\lor y \lor x \lor z~\lor \dots)=(\dots\lor y \lor z~\lor \dots).\]

\begin{example}
    Máme formuli
    \[\chi(x,y,z)=(x\lor y)\land y\land (z\lor \neg y)\land(\neg x\lor \neg y\lor \neg z).\]
    Aplikujme na ni proceduru \emph{Unit propagation}. Můžeme si všimnout,~že klauzule $y$ je jednotková a~proměnná $y$ tedy musí být nastavena na $1$.

    Protože $y=1$,~je zároveň splněna i~klauzule $x\lor y$,~takže ji můžeme odstranit. Zároveň odstraníme literál $\neg y$ ze zbývajících klauzulí $z\lor\neg y$ a~$\neg x\lor\neg y\lor\neg z$.
    Celkově dostaneme novou zjednodušenou formuli
    \[\chi^\prime(x,y,z)=z\land(\neg x\lor \neg z).\]
\end{example}

\subsubsection*{Pure literal elimination}

Ve vstupní formuli $\varphi$ může nastat situace,~kdy se zde určitá proměnná nachází pouze v~jedné polaritě. Tzn. pokud se ve formuli nachází proměnná $x$,~pak se v~ní nikde nenachází $\neg x$ a~naopak pokud se v~ní nachází $\neg x$,~pak v~ní nikde není $x$. V~obou případech můžeme rovnou určit hodnotu dané proměnné,~neboť tím splníme všechny klauzule,~v~nichž se vyskytuje. Ty můžeme odstranit,~neboť jsou tím pádem splněny.
\[\dots\land\overbrace{(\dots\lor \underbrace{x}_{1}\lor\dots)}^{1}\land\dots\land\overbrace{(\dots\lor \underbrace{x}_{1}\lor\dots)}^{1}\land\dots\]
Podobně pro negaci,~tj.
\[\dots\land\overbrace{(\dots\lor \neg\underbrace{x}_{0}\lor\dots)}^{1}\land\dots\land\overbrace{(\dots\lor \neg\underbrace{x}_{0}\lor\dots)}^{1}\land\dots\]

Poté,~co jsme si vysvětlili základní dvojici procedur,~se nyní nabízí trojice případů,~které je ještě třeba prozkoumat.
\begin{itemize}
    \item \emph{Co když postupným odstraňováním proměnných mi ve formuli vznikne prázdná klauzule?} Taková situace může nastat v~případě procedury Unit propagation,~kdy odstraňujeme proměnné v~opačné polaritě. Pokud nastane situace,~že je nějaká klauzule prázdná,~pak to znamená,~že daná klauzule pod zvoleným ohodnocením \textbf{není splnitelná} a~tudíž ani celá formule. V~takovou chvíli můžeme prohlásit formuli za \textbf{nesplnitelnou pod daným ohodnocením}\footnote{Je potřeba si uvědomit,~že ohodnocení proměnných si během výpočtu volíme,~tedy v~takovém případě to ještě nutně neznamená,~že je formule celkově nesplnitelná. Pouze víme,~že dané ohodnocení určitě není splňující}.
    \item \emph{Co když odstraníme postupně všechny klauzule?} K~takové situaci může dojít,~pokud jsou všechny klauzule splněny. V~takovou chvíli určitě víme,~že je původní formule \textbf{splnitelná}.
    \item \emph{Co když po aplikaci obou procedur stále nevíme,~zda je $\varphi$ splnitelná?} V~takovou chvíli nám nezbývá nic jiného než si zvolit nějakou dosud neohodnocenou proměnnou a~prozkoumat obě možnosti jejího ohodnocení. Vybereme tedy proměnnou $x$ a~zkusíme ji nastavit jednou na hodnotu $1$ a~podruhé na hodnotu $0$. Tím můžeme formuli $\varphi$ zjednodušit a~na nově vzniklou formuli $\varphi^\prime$ rekurzivně aplikovat stejný algoritmus. Pokud lze splnit zbytek formule $\varphi^\prime$,~pak lze splnit i~původní formuli $\varphi$.
    
    Toho můžeme dosáhnout tak,~že k~formuli $\varphi$ přidáme uměle jednotkovou klauzuli $x$,~resp. v~druhém případě $\neg x$. Tuto jednotkovou klauzuli si pak vybere procedura Unit propagation,~která proměnnou ohodnotí. Zavoláme tedy algoritmus DPLL rekurzivně na formule
    \[\varphi\land x\quad\text{a}\quad\varphi\land\neg x.\]
\end{itemize}

S~těmito informacemi můžeme dát dohromady pseudokód algoritmu.
\begin{algorithm}
    \caption{DPLL}
    \label{alg:dpll}
    \KwIn{Formule $\varphi$ v~CNF}
    \While{\textup{existuje jednotková klauzule $(\ell)$ ve formuli $\varphi$}}{
        $\varphi\gets\text{UnitPropagation}(\ell,~\varphi)$
    }
    \While{\textup{existuje literál $\ell$ v~jediné polaritě ve $\varphi$}}{
        $\varphi\gets\text{PureLiteralElimination}(\ell,\varphi)$
    }
    \tcp{Všechny klauzule jsou pravdivé}
    \lIf{je $\varphi$ prázdná}{\Return{$1$}}

    \tcp{Všechny literály v~klauzuli jsou nepravdivé}
    \lIf{\textup{obsahuje $\varphi$ prázdnou klauzuli}}{\Return{$0$}}
    \tcp{Zkusíme obě ohodnocení proměnné $x$}
    Vyber dosud neohodnocenou proměnnou $x$\\

    \Return{$\textup{DPLL}(\varphi\land x)\lor\textup{DPLL}(\varphi\land\neg x)$}

    \KwOut{$1$,~je-li formule $\varphi$ splnitelná,~jinak $0$.}
\end{algorithm}

Zkusme si algoritmus odsimulovat na příkladu.
\begin{example}
    Máme formuli
    \[\gamma(x_1,x_2,x_3,x_4)=(x_1\lor x_3)\land\neg x_4\land(\neg x_2\lor\neg x_3)\land x_1\land(\neg x_2\lor x_3)\land(\neg x_1\lor\neg x_3\lor x_3).\]
    Chceme pomocí algoritmu DPLL rozhodnout,~zda je splnitelná.

    Když se podíváme na pseudokód \ref{alg:dpll} výše,~tak první máme aplikovat proceduru \emph{Unit propagation}. Hledáme tedy jednotkové klauzule ve formuli $\gamma$,~dokud nějaké existují.
    \begin{itemize}
        \item První jednotková klauzule,~které si můžeme všimnout,~je $\neg x_4$. Hodnota proměnné $x_4$ tedy nutně musí být $0$. Klauzuli tedy odstraníme. Proměnná se nachází ve formuli pouze jednou,~tedy daný literál v~opačné polaritě nikde odstraňovat nemusíme. Máme nyní formuli
        \[\gamma^\prime(x_1,x_2,x_3,x_4)=(x_1\lor x_3)\land(\neg x_2\lor\neg x_3)\land x_1\land(\neg x_2\lor x_3)\land(\neg x_1\lor\neg x_3\lor x_3).\]
        \item Další jednotková klauzule je $x_1$. Proměnná $x_1$ musí být nutně $1$ a~klauzuli můžeme odstranit. Dále si však můžeme všimnout,~že tím pádem splněna i~klauzule $(x_1\lor x_3)$. Tu můžeme též odstranit. Literál $x_1$ se dále nachází v~opačné polaritě v~klauzuli $(\neg x_1\lor\neg x_3\lor x_3)$. Literál $\neg x_1$ můžeme také odstranit z~formule. Tím se nám formule $\gamma^\prime$ zjednoduší na
        \[\gamma^{\prime\prime}(x_1,x_2,x_3,x_4)=(\neg x_2\lor\neg x_3)\land (\neg x_2\lor x_3)\land(\neg x_3\lor x_3).\]
    \end{itemize}
    Nově vzniklá formule $\gamma^{\prime\prime}$ již žádné další jednotkové klauzule neobsahuje. Dále tedy můžeme aplikovat proceduru \emph{Pure literal elimination}.
    \begin{itemize}
        \item Formule $\gamma^{\prime\prime}$ obsahuje pouze jeden literál v~jediné polaritě,~a~to $\neg x_2$. Tím pádem hodnotu $x_2$ můžeme nastavit na $0$,~čímž splníme dvojici klauzulí $(\neg x_2\lor\neg x_3)$ a~$(\neg x_2\lor x_3)$. Ty můžeme odstranit,~čímž dostaneme formuli
        \[\gamma^{\prime\prime\prime}(x_1,x_2,x_3,x_4)=\neg x_3\lor x_3.\]
    \end{itemize}
    Protože formule $\gamma^{\prime\prime\prime}$ není ani prázdná,~ani neobsahuje prázdnou klauzuli,~vybereme náhodně jeden literál a~prozkoumáme jeho obě možná ohodnocení. V~tomto případě nám zbývají pouze literály $x_3$ a~$\neg x_3$. Vyberme např. literál $\neg x_3$. Nyní máme sestrojit dvojici nových formulí $(\neg x_3\lor x_3)\land\neg x_3$ a~$(\neg x_3\lor x_3)\land x_3$,~na které rekurzivně zavoláme DPLL.
    \begin{itemize}
        \item \textbf{Formule $(\neg x_3\lor x_3)\land\neg x_3$.} Opět začneme aplikací Unit propagation. K~původní formuli $\gamma^{\prime\prime\prime}$ jsme uměle přidali jednotkovou klauzuli $\neg x_3$. Z~ní je zjevné,~že $x_3$ musí být $0$. Tím zároveň splníme i~klauzuli $(\neg x_3\lor x_3)$,~kterou můžeme odstranit,~čímž dostáváme prázdnou formuli. Pure literal elimination aplikovat nemůžeme. Protože je nově vzniklá formule prázdná,~vrátíme na výstup $1$.
        \item \textbf{Formule $(\neg x_3\lor x_3)\land x_3$.} I~zde začneme eliminací jednotkových klauzulí. Tu zde představuje literál $x_3$,~jehož ohodnocení tedy musí být $1$. Tím opět splníme i~klauzuli $(\neg x_3\lor x_3)$,~jejímž odstraněním dostaneme zase prázdnou formuli. Tedy opět vrátíme $1$.
    \end{itemize}
    Obě ohodnocení pro $x_3$ jsou tedy splňující a~tedy i~původní formule $\gamma$ \emph{je splnitelná}.
\end{example}

Ukažme si ještě jeden jednodušší příklad.
\begin{example}
    Mějme formuli $\tau(x,y,z)=x\land (\neg x\lor y\lor z)\land \neg x$. Chceme opět pomocí DPLL určit splnitelnost.

    Jako první jednotkovou klauzuli můžeme vzít $x$. Tedy můžeme ji odstranit společně se všemi výskyty daného literálu v~opačné polaritě,~tj. $\neg x$. Tzn. dostaneme formuli
    \[\tau^\prime(x,y,z)=(y\lor z)\land().\]
    Všimněte si,~že odstraněním literálu $\neg x$ jsme dostali ve formuli prázdnou klauzuli\footnote{Z~jednotkové klauzule jsme odebrali její jediný literál v~rámci procedury Unit propagation; klauzule samotná ve formuli zůstala.} $()$. Aplikací Pure literal elimination se zbavíme i~literálů $y$ a~$z$.
    \[\tau^{\prime\prime}(x,y,z)=()\]
    Formule $\tau^{\prime\prime}$ obsahuje pouze prázdnou klauzuli. V~tuto chvíli algoritmus prohlásí formuli $\tau$ za \emph{nesplnitelnou}.
\end{example}

Poslední otázka,~která se nabízí,~je (jako obvykle) časová složitost algoritmu. Vezmeme-li v~potaz nejhorší případ,~časová složitost bohužel nevychází úplně příznivě.
\begin{theorem}[Složitost DPLL]
    Nechť formule $\varphi$ obsahuje $v$ proměnných a~její délka je $s$. Jednoduchá implementace DPLL doběhne v~čase $\bigO(s2^v)$. Při kopírování formule v~každé úrovni rekurze spotřebuje nejvýše $\bigO(sv)$ pomocné paměti.
\end{theorem}
\begin{proof}
    V~nejhorším případě nepomůže ani jednotková propagace,~ani eliminace čistých literálů. Algoritmus pak zvolí proměnnou a~vytvoří dvě větve,~jednu pro hodnotu $0$ a~druhou pro hodnotu $1$. Hloubka stromu je nejvýše $v$ a~počet jeho uzlů nejvýše
    \[1+2+4+\dots+2^v=2^{v+1}-1=\bigO(2^v).\]
    V~každém uzlu lze naivně projít celou formuli v~čase $\bigO(s)$,~odkud plyne čas $\bigO(s2^v)$. Rekurze probíhá do hloubky,~takže současně uchováváme nejvýše $v$ kopií formule délky nejvýše $s$; to dává $\bigO(sv)$ pomocné paměti.
\end{proof}
\notebox{Praktické \ensuremath{\SAT} solvery formuli při každém volání celou nekopírují. Změny zapisují do pomocných struktur a~při návratu je vracejí zpět,~takže paměťový odhad lze podstatně zlepšit.}

S~exponenciální časovou složitostí se sotva spokojíme,~avšak v~tomto případě je dobré si říci,~že zkoumáme složitost v~nejhorším případě a~v~praxi bývají zkoumané logické formule poměrně rychle řešitelné pomocí DPLL.
\importantbox{Problém je však dodnes ten,~že na rozhodnutí splnitelnosti logické formule (i~těch v~CNF) není známý žádný algoritmus,~který by běžel v~polynomiálním čase,~tj. $\bigO(n^k)$,~pro nějaké $k$. K~tomuto tématu se ještě vrátíme v~sekci \ref{sec:p_np}.}

\section{Problém 3-SAT}\label{sec:3sat}

Zkusme nyní původní problém ještě zjednodušit. Přidáme podmínku, že každá klauzule může obsahovat nejvýše 3 literály. Tím dostáváme variantu problému SAT známý jako \emph{3-SAT}. O formuli, která splňuje zmíněnou formu, budeme říkat, že je v 3-CNF.

\problem{Splnitelnost logické formule v 3-CNF (3-SAT)}{Formule $\varphi$ v 3-CNF}{1, je-li $\varphi$ splnitelná, jinak $0$.}

Nabízí se otázka, jestli jsme si skutečně nějak pomohli. Zdánlivě ano, neboť máme kratší klauzule, avšak 3-SAT převdstavuje, překvapivě, stejně těžkou úlohu jako původní problém SAT. Proč? Lze totiž ukázat, že pokud bychom uměli rychle řešit 3-SAT, uměli bychom rychle řešit i původní SAT a naopak. To ukážeme tak, že jeden problém v jistém smyslu ``převést'' na ten druhý. Pro určitou instanci (tedy formuli) problému SAT můžeme sestrojit odpovídající instanci problému 3-SAT, jejímž vyřešením získáme odpověď i na původní problém. Podívejme se na tuto situaci blíže v podsekci

\subsection{Převod 3-SAT na SAT}\label{subsec:prevod_sat_3sat}

Popíšeme si nyní obecný způsob, jak převést libovolnou formuli v CNF (instance problému SAT) na formuli v 3-CNF (instance problému 3-SAT). Předpokládejme, že máme zadanou nějakou formuli $\varphi$ v CNF a chceme pro ni sestrojit formuli $\psi$ v 3-CNF, takovou, aby byla splnitelná právě tehdy, když je splnitelná i původní formule $\varphi$. Pro formuli $\varphi$ mohou nastat 2 případy:
\begin{enumerate}
    \item formule obsahuje klauzule délky $3$ a kratší,
    \item nebo ve formuli existuje klauzule délky $k>3$.
\end{enumerate}
V prvním případě není co řešit, neboť $\varphi$ již je v 3-CNF, tedy stačí položit $\psi=\varphi$ a jsme hotovi. V druhém případě musíme dlouhou klauzuli o $k$ literálech rozložit na kratší. Tuto klauzuli\footnote{Pokud $\varphi$ obsahuje více dlouhých klauzulí, postup opakujeme pro každou z nich zvlášť.} rozložíme na dvojici nových klauzulí pomocí nové proměnné $x$. Označme si literály dané dlouhé klauzule $\ell_1,\ell_2,\dots\ell_k$. První dvojici $\ell_1$ a $\ell_2$ si označíme $\alpha$ a zbytek $\beta$.
\[\underbrace{\ell_1\lor\ell_2}_{\alpha}\lor\overbrace{\ell_3\lor\dots\lor\ell_k}^{\beta}=\alpha\lor\beta\]
Nyní přidáme novou proměnnou $x$ a klauzuli $\alpha\lor\beta$ rozdělíme na dvojici klauzulí následovně:
\[(\alpha\lor x)\land(\beta\lor\neg x).\]
Podívejme se nyní na délky nových klauzulí $\alpha\lor x$ a $\beta\lor\neg x$.
\begin{itemize}
    \item Klauzule $\alpha\lor x=\ell_1\lor\ell_2\lor x$ obsahuje $3$ literály, takže s ní dále nic provádět nemusíme.
    \item Klauzule $\beta\lor\neg x=\ell_3\lor\dots\lor\ell_k\lor\neg x$ obsahuje $k-1$ literálů. To znamená, že jsme zkrátili původní klazuli o jeden literál. Pokud $k-1>3$, můžeme postup pro tuto klauzuli opakovat.
\end{itemize}
Nyní zbývá dořešit, jak to bude s hodnotou proměnné $x$. Podobně jako u Tseitinovy transformace (viz podsekce \ref{subsec:cnf_tseitin}), i zde bude její hodnota záviset na hodnotách ostatních proměnných, konkrétně výrazu $\alpha$.
\begin{itemize}
    \item Pokud $\alpha=1$, pak nastavíme $x=0$. Původní klauzule $\ell_1\lor\dots\lor\ell_k=\alpha\lor\beta$ je splněna díky $\alpha$. Nová dvojice klauzulí je pak splněna pomocí $\alpha$ a $x$:
    \[(\alpha\lor x)\land(\beta\lor\neg x)=(1\lor 0)\land(\beta\lor 1)=1\land 1=1.\]
    \item Pokud $\alpha=0$, pak nastavíme $x=1$. Klauzule $\alpha\lor\beta$ musí být splněna díky $\beta$. Potom platí:
    \[(\alpha\lor x)\land(\beta\lor\neg x)=(0\lor 1)\land(1\lor 0)=1\land 1=1.\]
    Pokud by však $\beta=0$, pak by původní klauzule $\alpha\lor\beta$ nebyla splněna a potom
    \[(\alpha\lor x)\land(\beta\lor\neg x)=(0\lor 1)\land(0\lor 0)=1\land 0=0,\]
    tzn. ani nová dvojice klauzulí by nebyla splněna, což chceme.
\end{itemize}
Tím jsme dostali obecný postup pro převod formule z CNF do 3-CNF.
\begin{example}
    Máme formuli $\varphi(a,b,c,d)=a\lor b\lor \neg c\lor \neg b\lor d\lor \neg d$ o jedné klauzuli.
\end{example}
\section{Problém nezávislé množiny v grafu (NzMna)}\label{sec:nzmna}
\section{Problém kliky v grafu}\label{sec:klika}
\section{Třídy problémů \P{} a \NP}\label{sec:p_np}
